\documentclass[11pt]{article}
\usepackage[margin=1in]{geometry}
\usepackage{amsmath,amssymb,amsthm}
\usepackage{hyperref}
\hypersetup{hidelinks}

\newtheorem{theorem}{Theorem}
\newtheorem{lemma}[theorem]{Lemma}
\newtheorem{corollary}[theorem]{Corollary}
\theoremstyle{definition}
\newtheorem{definition}{Definition}
\theoremstyle{remark}
\newtheorem{remark}{Remark}

\setlength{\parskip}{4pt plus 1pt minus 1pt}
\setlength{\parindent}{0pt}

\title{Agency and Learning Efficiency as the\\Structural Invariants of Intelligence}
\author{Patrick D.\ McCarthy}
\date{}

\begin{document}
\maketitle

\begin{abstract}
We prove that for any entity with structure $(K, A, \sigma, \gamma)$
persisting under bounded active context, survival pressure, and
positive novelty rate $\nu(E) > 0$: (1)~agency $\gamma > 0$ is
necessary: the lineage argument shows that $\gamma = 0$ lineages contract
and are eliminated; (2)~learning efficiency $\eta_M > 0$ is necessary:
without it, the active context $L_n$ fills without relief and
inference degrades to failure; (3)~the action space $A$ (informed
actions) and the higher-order function space $M$ (learning) are
type-irreducible: no learning process changes a function's type
signature; (4)~every other property of the entity can equal zero for
some persistent entity; (5)~therefore $I(E) = \gamma \cdot \eta_M$ is the
unique measure of intelligence, derived not stipulated. This completes
the structural theory: from bounded context alone, the complete
architecture of intelligence follows as necessary consequences.
\end{abstract}

% ------------------------------------------------------------------
\section{Introduction}
\label{sec:intro}
% ------------------------------------------------------------------

This paper completes a chain of results deriving the architecture of
intelligence from a single generative constraint: bounded active context.

Paper~1 \cite{McCarthy2026a} established that any entity accumulating
propositions under bounded context necessarily maintains a directed
graph with typed edges. Paper~2 \cite{McCarthy2026b} established that
bounded context under growing action count requires escalation
to enable top-down allocation. Paper~3 \cite{McCarthy2026c}
established that survival pressure makes knowledge and action
inseparable, with an asymmetric retention criterion.
Paper~4 \cite{McCarthy2026d}
established that agency with feedback from the world produces gradient
descent on uncertainty, with convergence, and that cost-of-failure
generates an encoding hierarchy with permanence.

One question remains: \textbf{what properties must every persisting
entity have?} We prove that the answer is exactly two: agency $\gamma > 0$
and learning efficiency $\eta_M > 0$. Their product is intelligence.

This result engages a long-standing difficulty. Dupoux, LeCun, and Malik
\cite{DupouxLeCunMalik2026} deliberately avoid defining intelligence,
instead characterizing it as the capacity for autonomous learning. We
derive the definition they avoided giving. Hutter \cite{Hutter2005} and
Legg and Hutter \cite{LeggHutter2007} define intelligence as average
performance over computable environments; we define it as rate of
improvement, independent of current performance level. Chollet
\cite{Chollet2019} measures skill-acquisition efficiency, structurally
close to $\eta_M$ but treated as a benchmark metric rather than a
structural invariant. Schmidhuber \cite{Schmidhuber2010} identifies
curiosity as compression progress, which corresponds to $\eta_M$, but
treats drive as a design objective rather than a necessary condition of
persistence. Friston's free energy principle \cite{Friston2010} derives
minimization behavior from self-organization but does not establish the
minimization \emph{imperative}: why $\gamma > 0$ is structurally necessary.

The contribution of this paper is to prove, via a selection argument on
lineages, that $\gamma > 0$ and $\eta_M > 0$ are the unique structural
invariants of persistence, and that $I(E) = \gamma \cdot \eta_M$ is the only
consistent intelligence measure.

% ------------------------------------------------------------------
\section{Formal Framework}
\label{sec:framework}
% ------------------------------------------------------------------

We build on the definitions of Papers~1--4 and introduce the structures
needed for the present results. Throughout, $W$ is a measurable space of
world states, $K$ is a knowledge base of
propositions about $W$, $C_n$ is the bounded active context capacity,
$A$ is the action space (actions informed by knowledge), and $U(w,K) =
-\log P(w \mid K)$ is the uncertainty of $K$ about state $w$
\cite{Shannon1948, Cover2006}.

\begin{definition}[Novelty Rate]
\label{def:novelty}
The \emph{novelty rate} $\nu(E)$ of entity $E$ in world $W$ is the rate
at which $E$ encounters world states $w$ where $U(w, K) > 0$:
\[
\nu(E) = \lim_{T \to \infty} \frac{1}{T}
\bigl|\{t \leq T : U(w_t, K(t)) > 0\}\bigr|
\]
Novelty rate is a property of the gap between $W$ and $K$, not of $W$
alone. A world with high entropy but an entity whose $K$ covers all
reachable states has $\nu = 0$. A world with low entropy but an entity
with impoverished $K$ has $\nu > 0$. The distinction matters: entropy is
a property of the environment; novelty is a property of the
entity--environment coupling.
\end{definition}

\begin{definition}[Free Context]
\label{def:free}
The \emph{free context} of entity $E$ at time $t$ is the active context
capacity not consumed by unresolved content:
\[
C_{\mathrm{free}}(t) = C_n - |\mathrm{ctx}_{\mathrm{unresolved}}(t)|
\]
where $\mathrm{ctx}_{\mathrm{unresolved}}(t)$ is the set of
propositions in active context that have not yet been encoded at a lower
level or resolved by $M$. Free context is the entity's reserve capacity
for handling novel states. Bounding uncertainty on known problems frees
context; novel problems consume it. The architecture derived in
Papers~1--4, graph structure, escalation, encoding permanence, exists
to maximize $C_{\mathrm{free}}$: each result describes a mechanism that
moves resolved content out of $L_n$, preserving capacity for what the
entity has not yet encountered.
\end{definition}

\begin{definition}[Higher-Order Function Space $M$]
\label{def:M}
Let $A : K \times W_{\mathrm{obs}} \to \Delta \mathcal{A}$ be the
action space of informed actions, where $W_{\mathrm{obs}}$ denotes
observable world states and $\Delta \mathcal{A}$ the simplex over
available actions. The \emph{higher-order function space} is
\[
M : (A, K) \times L \;\to\; (A, K)
\]
where $L$ is a learning signal (feedback from $W$ following action).
$M$ is categorically different from $A$: elements of $A$ map knowledge
and observations to action distributions; elements of $M$ map
action--knowledge pairs and learning signals to updated
action--knowledge pairs.

$M$ contains whatever adaptive capacity has not yet been committed to
specific knowledge or action. As the entity learns, components of $M$
resolve into concrete elements of $K$ or $A$---the same normalization
by which shared propositions precipitate from action into knowledge
\cite{McCarthy2026c}. What remains in $M$ is the entity's capacity
to produce further such resolutions.

The hierarchy is open upward: $M^{(2)}$ operates on $M^{(1)}$, yielding
meta-learning; $M^{(n+1)}$ operates on $M^{(n)}$. No natural
termination point is established (see Section~\ref{sec:open}).
\end{definition}

\begin{definition}[$M$ Efficiency]
\label{def:eta}
Let $G(a, K) = I(p_a; W \mid K)$ denote the mutual information between
the feedback from action $a$ and the true world state $W$, conditional
on knowledge $K$---the \emph{grounding} of $a$. Given a learning event
where $(A', K') = M((A, K), L)$, define the \emph{learning efficiency}:
\[
\eta_M \;=\; \mathbb{E}_{L \mid \text{engaged}}\!\left[
\frac{\max_{a' \in A'} G(a', K') - \max_{a \in A} G(a, K)}{|L|}
\right]
\]
The expectation is taken over learning signals $L$ \emph{conditional on
the entity engaging $M$} (i.e., conditional on the gate $\gamma$
opening). This conditioning is essential: $\gamma$ may be directed
(the entity may selectively engage on high-quality signals), so
$\mathbb{E}[{\cdot} \mid \text{engaged}]$ may differ from the
unconditional $\mathbb{E}_L[{\cdot}]$. Conditioning on engagement
ensures that the product $\gamma \cdot \eta_M$ is the expected
improvement rate by the law of total expectation, regardless of
whether $\gamma$ is content-blind or directed.

Both numerator and denominator are measured in bits; $\eta_M$ is
dimensionless. The three regimes:
\begin{itemize}
\item $\eta_M > 0$: the entity learns: grounding improves per unit
      of learning signal, conditional on engagement.
\item $\eta_M = 0$: engagement produces no improvement: $M$ is
      engaged but yields no grounding gain.
\item $\eta_M < 0$ temporarily: exploration, creative destruction,
      locally irrational steps that may increase long-run survival.
      Persistently $\eta_M < 0$: the entity does not persist.
\end{itemize}
\end{definition}

\begin{remark}[Measurement]
\label{rem:measurement}
Computing $\eta_M$ requires evaluating $G(a, K) = I(p_a; W \mid K)$,
which depends on the true joint distribution $P(p_a, W, K)$. The entity
cannot compute its own $\eta_M$; it would need the very information
it is trying to learn. This is a structural feature, not a defect:
self-knowledge of learning rate would require solving the problem that
learning is meant to solve. Both $\gamma$ and $\eta_M$ are measurable by an
external observer with access to the entity's trajectory and outcomes.
\end{remark}

\begin{definition}[Intelligence]
\label{def:intelligence}
For an entity $E$ with agency $\gamma$ and learning efficiency $\eta_M$:
\[
I(E) = \gamma \cdot \eta_M
\]
This is stated here as a definition. Theorem~\ref{thm:invariants}
establishes that it is the unique measure consistent with the structural
invariants of persistence.
\end{definition}

\begin{definition}[Agency]
\label{def:agency}
Agency $\gamma \in [0, 1]$ is the probability that the entity engages its
adaptive capacity $M$ in response to a learning signal. Equivalently,
$\gamma$ is the response gate: the fraction of environmental signals that
trigger an update to $(A, K)$.

The operational signature of $\gamma > 0$ is that the entity bounds
uncertainty beyond what immediate survival requires, freeing context
capacity for novel states. An entity with $\gamma = 0$ responds to no
signal and accumulates no free context. An entity with $\gamma > 0$
engages $M$ on signals, resolves uncertainty, promotes actions to lower
encoding levels, and thereby maintains $C_{\mathrm{free}} > 0$
(Definition~\ref{def:free}).

The distinction between $\gamma = 0$ and $\gamma > 0$ is not between
entities that face threats and those that do not. It is between entities
that build reserve capacity for states they have not yet encountered and
those that do not.
\end{definition}

% ------------------------------------------------------------------
\section{Type Irreducibility of $A$ and $M$}
\label{sec:type}
% ------------------------------------------------------------------

\begin{theorem}[Type Irreducibility of $A$ and $M$]
\label{thm:type}
No element of $A$ can become an element of $M$ through any learning
process. No element of $M$ can become an element of $A$ through any
learning process. The two spaces are type-irreducible.
\end{theorem}

\begin{proof}
The type signatures are distinct:
\begin{align*}
A &: K \times W_{\mathrm{obs}} \to \Delta \mathcal{A} \\
M &: (A, K) \times L \to (A, K)
\end{align*}
$A$ takes knowledge and observations as input and returns an action
distribution. $M$ takes an action--knowledge pair and a learning
signal as input and returns an updated action--knowledge pair. The
domains and codomains are different types.

No learning process changes a function's type signature. An element
$a \in A$ that achieves certainty (its output concentrates on a single
action for all inputs) promotes to a lower encoding level in the
hierarchy established in \cite{McCarthy2026d}, but this is a change
of encoding level, not a change of type. The promoted $a$ still maps
$K \times W_{\mathrm{obs}} \to \Delta \mathcal{A}$; it is still an
element of $A$.

Conversely, $M$ takes $(A, K)$ as input. An element of $A$ does not
accept $(A, K)$ as input; it accepts $K \times W_{\mathrm{obs}}$.
No amount of learning can change what a function accepts as input. The
type distinction is categorical, not parametric.

\emph{Why unification fails.} A natural objection holds that the type
distinction is definitional: we defined $A$ and $M$ to have different
signatures, so of course they are different. The non-trivial content
is that any entity attempting to collapse $A$ and $M$ into a single
type $T$ faces an impossibility. Suppose $T : K \times W_{\mathrm{obs}}
\times (T, K) \times L \to \Delta \mathcal{A} \times (T, K)$. An
element of $T$ must simultaneously accept observations (to select
actions) and accept its own space as input (to modify itself). But
self-application $T(T)$ is ill-typed in the simply-typed lambda
calculus \cite{Barendregt1984}. Polymorphic systems (System~F, the
Calculus of Constructions) permit forms of self-application via
universal quantification, but the quantified function does not modify
its own type---it operates uniformly across types. The relevant
constraint is not syntactic self-application but \emph{self-modification
at the same meta-level}: a function that rewrites its own
implementation while executing. This requires stratification---exactly
the $A$/$M$ separation.

An anticipated objection: programs routinely modify themselves
(\texttt{eval}, metaclasses, Lisp macros). But in every such case,
the code being modified operates at a lower meta-level than the code
doing the modification. The interpreter that evaluates \texttt{eval}
is not the code being evaluated. This is the $A$/$M$ separation
implemented in software: the runtime ($M$) modifies the program ($A$),
and these occupy different levels of the execution hierarchy. Languages
that appear to collapse the levels achieve it by implicit
stratification, not by genuine self-application. The irreducibility
is not our assumption but a consequence of type-theoretic consistency.
\end{proof}

\begin{corollary}[Distinct Roles]
\label{cor:roles}
The type irreducibility establishes distinct functional roles:
\begin{itemize}
\item $A$ navigates $W$ given $K$: selecting actions in the current
      world state using current knowledge.
\item $M$ evolves $(A, K)$: modifying the entity's action repertoire
      and knowledge base in response to learning signals.
\end{itemize}
Concrete instances of $M$ acting on different targets:
\begin{itemize}
\item \textbf{Neuroplasticity}: $M$ acting on $A$: modifying which
      actions are available and how they respond to inputs.
\item \textbf{Insight}: $M$ acting on $K$: restructuring the
      knowledge base to incorporate a new proposition that changes the
      graph topology.
\item \textbf{Habit formation}: $M$ promoting from $L_n$ to lower
      encoding levels, the process described in
      \cite{McCarthy2026d}, Theorem~4.
\item \textbf{Evolution}: $M$ acting on $M$ itself, the highest
      level of the hierarchy, operating at the most stringent rigor
      threshold.
\end{itemize}
\end{corollary}

% ------------------------------------------------------------------
\section{Necessity of $\gamma$ and $\eta_M$}
\label{sec:necessity}
% ------------------------------------------------------------------

\begin{definition}[Persistence]
\label{def:persistence}
An entity $E$ \emph{persists} if it survives for all $t$: there is no
finite $T$ at which $U(w_T, K(T)) > U_{\mathrm{lethal}}^d$ terminates
the entity. Persistence is indefinite survival, not survival to some
fixed horizon.
\end{definition}

\begin{theorem}[Structural Invariants of Persistence]
\label{thm:invariants}
For any entity $E$ persisting (Definition~\ref{def:persistence}) under
bounded active context $C_n$, survival threshold
$U_{\mathrm{lethal}}^d$, with positive novelty rate $\nu(E) > 0$
(Definition~\ref{def:novelty}):
\begin{enumerate}
\item[(i)] $\gamma(E) > 0$ necessarily.
\item[(ii)] $\eta_M(E) > 0$ necessarily (in expectation over the
      entity's trajectory).
\item[(iii)] Every other scalar property of $E$ except $\sigma > 0$
      (sensing) can equal zero for some persistent entity.
\item[(iv)] Therefore $I(E) = \gamma \cdot \eta_M$ is the natural measure
      of intelligence, identified by what survival requires.
\end{enumerate}

When $\nu(E) = 0$, the entity faces no novel states and $\gamma = 0$ is
viable: no free context is needed because nothing unexpected arrives.
Such entities persist without intelligence. Intelligence is what positive
novelty rate selects for.
\end{theorem}

\textbf{Scope note.} This theorem establishes $\gamma > 0$ and $\eta_M > 0$
as \emph{necessary} conditions for persistence. The proof is a selection
argument: entities lacking either property are eliminated. The theorem
does \emph{not} explain the \emph{origin} of $\gamma$: how $\gamma > 0$
structures bootstrap from $\gamma = 0$ substrates is an open question (see
Section~\ref{sec:open}).

\textbf{Scope clarification.} The theorem applies to entities with the
structure $E = (K, A, \sigma, \gamma)$. This is a presupposition, not a
conclusion: the theorem does not derive the \emph{existence} of agency
from non-agentive precursors. What it establishes is an
\emph{elimination result}: among entities that have this structure,
those with $\gamma = 0$ do not persist under $\nu > 0$. The
contribution is identifying \emph{why} $\gamma = 0$ is untenable---the
free-context argument---not explaining how $\gamma > 0$ arises from
$\gamma = 0$ substrates (which remains an open question;
Section~\ref{sec:open}). Inert objects (rocks, tables) lack this
structure entirely and are outside the theorem's domain.

\begin{proof}

\textbf{Part 1: Necessity of $\gamma > 0$: The Free Context Argument.}

Suppose $\gamma(E) = 0$. By Definition~\ref{def:agency}, the entity never
engages $M$: its adaptive capacity is zero. Therefore $(A(t), K(t)) =
(A(0), K(0))$ for all $t$; the entity's action--knowledge pair is
fixed at its initial state.

The entity has positive novelty rate $\nu(E) > 0$
(Definition~\ref{def:novelty}): it encounters world states where
$U(w, K) > 0$ at a positive rate. Each novel state consumes active
context capacity. With $\gamma = 0$, no learning occurs: no uncertainty
is resolved, no action is promoted to a lower encoding level, no
context is freed. $C_{\mathrm{free}}(t) \to 0$ monotonically as
unresolved content accumulates. When $C_{\mathrm{free}} = 0$, the
entity cannot load the relevant subgraph $K^{[a]}$ for any new problem
$a$: inference quality $\rho \to 0$. The entity acts without relevant
knowledge. It encounters $w$ where $U(w, K(0)) > U_{\mathrm{lethal}}^d$
and does not survive.

The argument is self-contained at the individual level and requires no
competition. For any initial $K(0)$, because $\nu > 0$ the entity
encounters novel states at positive rate. Each novel state consumes
context capacity that is never freed: freeing context requires either
promoting actions to lower encoding levels or resolving uncertainty,
both of which require engaging $M$ \cite{McCarthy2026b, McCarthy2026d}.
With $\gamma = 0$, no such engagement occurs. Therefore there exists a
finite time $T^*$ at which $C_{\mathrm{free}}(T^*) = 0$. After $T^*$,
the entity cannot process any novel state. Since $\nu > 0$ continues
to present states where $U(w, K(0)) > 0$, the entity eventually
encounters $w$ with $U(w, K(0)) > U_{\mathrm{lethal}}^d$ and does not
survive. No competitor is required: the environment alone is
sufficient.

The argument extends to other scales:
\begin{itemize}
\item \textbf{Lineage}: a $\gamma = 0$ lineage disperses copies with
      fixed $(A(0), K(0))$ across $W$. Each copy faces the individual
      argument above; the lineage contracts as copies fail.
\item \textbf{Institutional}: organizations with $\gamma = 0$ at the
      organizational level face the same free-context exhaustion.
\end{itemize}

The critical distinction is between $\nu = 0$ and $\nu > 0$. An entity
in a zero-novelty environment (a deep-sea vent organism whose $K$
covers all reachable states, a thermostat in a stable room) faces no
context pressure and persists with $\gamma = 0$. The framework does not
predict it dies. It predicts it does not need intelligence. The
necessity of $\gamma > 0$ is conditional on the environment presenting
states the entity has not yet bounded.

Positive entropy in $W$ is the upstream cause of positive novelty rate:
energy gradients produce states not contained in any finite $K$
\cite{Cover2006, Ashby1956}. But entropy is a property of the
environment; novelty rate is a property of the entity--environment
coupling. Entropy guarantees that $\nu > 0$ eventually for any entity
whose niche is not permanently isolated from the rest of $W$. The
``eventually'' matters: a $\gamma = 0$ entity in a temporarily stable
niche persists until entropy-driven change reaches it. The theorem's
premise $\nu > 0$ is the condition that change has arrived.

Therefore $\gamma(E) > 0$ is necessary for persistence under $\nu > 0$.

\medskip
\textbf{Part 2: Necessity of $\eta_M > 0$.}

Suppose $\gamma > 0$ but $\eta_M \leq 0$.

\emph{Case $\eta_M = 0$}: The entity has drive ($\gamma > 0$) but $M$
produces no improvement in grounding, conditional on engagement.
$(A, K)$ is effectively fixed despite engagement. As in Part~1,
the free-context exhaustion argument applies: novel states consume
context that is never freed by promotion, and $\rho \to 0$.

As with $\gamma$, the exception is $\nu = 0$: an entity whose $K(0)$
already covers all reachable states faces no novel demands and persists
with $\eta_M = 0$. The necessity of $\eta_M > 0$ is conditional on
$\nu > 0$, exactly as for $\gamma$.

Additionally, $C_n$ is bounded. Novel world states generate learning
signals that fill the active context $L_n$ with unresolved content. By
the escalation result \cite{McCarthy2026b}, inference quality $\rho$
degrades as $L_n$ fills without promotion to lower encoding levels.
Promotion requires $\eta_M > 0$: achieving the rigor threshold
$\theta_i$ at encoding level $i$ requires $M$ to improve the relevant
action to the required certainty (Theorem~4 of \cite{McCarthy2026d}).
With $\eta_M = 0$, no action reaches threshold, $L_n$ is never
relieved, and $\rho \to 0$. Inference fails.

\emph{Case $\eta_M < 0$ persistently}: $(A, K)$ worsens in expectation
with each learning event. Grounding decreases; uncertainty increases.
This accelerates the failure mode of the $\eta_M = 0$ case.

The persistence condition is:
\[
\mathbb{E}_t[\eta_M] > 0 \quad \text{over the entity's trajectory}
\]
not $\eta_M(t) > 0$ at every time step. Temporary $\eta_M < 0$
(exploration, creative destruction) is compatible with persistence and
can increase long-run survival. What is required is that the
trajectory-averaged learning efficiency is positive.

Therefore $\mathbb{E}_t[\eta_M] > 0$ is necessary for persistence.

\medskip
\textbf{Part 3: Minimality: every other property can equal zero.}

For every other scalar property of $E$, there exists a persistent
entity for which that property equals zero:
\begin{itemize}
\item $K(0) = \varnothing$: an entity with $\gamma > 0$ and $\eta_M > 0$
      survives starting from empty knowledge; it builds $K$ from
      monotone descent on $H(W \mid K)$ \cite{McCarthy2026d}.
\item $A(0) = \{\mathrm{id}\}$: the identity action still returns
      feedback from $W$, grounding the first proposition.
\item $C_n = C_{\min}$: minimal context capacity. Graph structure
      \cite{McCarthy2026a} maintains $\rho = 1$ by retrieving exactly
      the relevant subgraph per problem.
\item Substrate, observable window $W_{\mathrm{obs}}$: these are free
      parameters. They constrain the rate and form of learning but not
      its possibility.
\end{itemize}

The exception is sensing: $\sigma > 0$ is necessary because $\gamma > 0$
requires learning signals, and learning signals require sensing. An
entity with $\sigma = 0$ receives no observations from $W$, $M$
receives no input, and no update occurs regardless of $\gamma$. However,
$\sigma > 0$ is a precondition for the entity to be an entity in the
relevant sense: without sensing, there is no $W_{\mathrm{obs}}$, no
feedback, and no action--world coupling. The theorem's domain
(Definition~\ref{def:agency}) presupposes $\sigma > 0$ as part of
entity structure, not as a separate invariant.

No other scalar property has the survival-critical necessity established
in Parts~1 and~2.

\medskip
\textbf{Part 4: $I(E) = \gamma \cdot \eta_M$ follows by arithmetic.}

Parts~1--3 are the substantive content: $\gamma$ and $\eta_M$ are the
unique structural invariants of persistence. Given that $\gamma$ is a
probability and $\eta_M$ a conditional rate, their product is the
expected improvement rate by the law of total expectation
(Corollary~\ref{cor:measure}). The discovery is in Parts~1--3; the
product form is a consequence of what the invariants \emph{are}, not
a separate result.
\end{proof}

% ------------------------------------------------------------------
\section{$I(E) = \gamma \cdot \eta_M$ Is Structurally Necessary}
\label{sec:measure}
% ------------------------------------------------------------------

\begin{corollary}[The Intelligence Measure]
\label{cor:measure}
$\gamma$ and $\eta_M$ are the unique structural invariants of persistence
(Theorem~\ref{thm:invariants}). The expected improvement rate of the
entity is:
\[
I(E) = \gamma \cdot \eta_M
\]
This is not selected from a family of candidates but derived from the
types of the invariants: $\gamma$ is a probability; $\eta_M$ is a
conditional rate. Their product is the expected rate by the law of
total expectation.
\end{corollary}

\begin{proof}
$\gamma \in [0, 1]$ is the marginal probability of engaging $M$.
$\eta_M$ is the expected improvement rate \emph{conditional on
engagement} (Definition~\ref{def:eta}). By the law of total
expectation:
\[
\mathbb{E}[\text{improvement rate}]
= \gamma \cdot \mathbb{E}[\text{rate} \mid \text{engaged}]
+ (1 - \gamma) \cdot 0
= \gamma \cdot \eta_M
\]
This identity holds whether $\gamma$ is content-blind (a fixed coin
flip) or directed (the entity selectively engages on promising
signals). The conditioning in Definition~\ref{def:eta} absorbs the
selection effect: $\eta_M$ already reflects whichever signals the gate
admits. The product is exact, not an approximation.

Many functions $g(\gamma, \eta_M)$ are zero on the axes and positive in
the interior ($\gamma^2 \cdot \eta_M$, $\gamma \cdot \eta_M^2$, etc.).
What distinguishes the product is not a uniqueness argument over
function space but the \emph{types of the operands}. $\gamma$ is a
probability of engagement; $\eta_M$ is the rate conditional on
engagement. Total expectation gives the unconditional rate. Any other
combination would require interpreting $\gamma$ or $\eta_M$ as something
other than what they are. The measure is derived from the structure of
the invariants, not stipulated.

The product $\gamma \cdot \eta_M$ has the same structure as the retention
criterion of Theorem~1 in \cite{McCarthy2026c}: probability of needing
it $\times$ conditional value given need. Both are instances of
expected value under a binary gate.
\end{proof}

\begin{remark}[Exploration and Creative Destruction]
\label{rem:exploration}
Temporarily negative $\eta_M$, locally irrational steps that reduce
grounding in the short run, is survival-optimal when the persistence
condition requires only $\mathbb{E}_t[\eta_M] > 0$ over the trajectory,
not $\eta_M \geq 0$ at every step. Exploration (trying actions with
uncertain outcomes) and creative destruction (dismantling existing
$(A, K)$ structure to rebuild it better) are the signature of $M$
sophisticated enough to optimize over trajectories rather than
individual steps. An entity restricted to $\eta_M \geq 0$ at every step
is trapped in local optima; an entity that can tolerate temporary
$\eta_M < 0$ accesses a larger region of the $(A, K)$ landscape.
\end{remark}

% ------------------------------------------------------------------
\section{Engagement with Prior Work}
\label{sec:prior}
% ------------------------------------------------------------------

\textbf{Dupoux, LeCun, and Malik \cite{DupouxLeCunMalik2026}.}
They deliberately avoid defining intelligence, instead characterizing
it as the capacity for autonomous learning. Theorem~\ref{thm:invariants}
establishes that this is structurally correct: $\gamma > 0$ (autonomous
engagement) and $\eta_M > 0$ (learning efficiency) are the unique
invariants. Their product is the definition they avoided giving. Their
dual-system architecture (observation + action with meta-control routing)
corresponds to the $A$/$M$ distinction established in
Theorem~\ref{thm:type}, with meta-control corresponding to $M^{(2)}$.

\textbf{Hutter \cite{Hutter2005} and Legg--Hutter \cite{LeggHutter2007}.}
They define intelligence as average performance over computable
environments, formalized via AIXI. Three structural differences from the
present account: (1)~AIXI assumes unbounded computation; we have bounded
$C_n$. (2)~They define intelligence as performance level; we define it as
rate of improvement ($\eta_M$), independent of current performance. A
novice chess player learning rapidly has higher $\eta_M$ than a grandmaster
at plateau, despite lower performance. (3)~Their environment class is all
computable functions; ours is physical $W$ with positive entropy, a
narrower but physically grounded class.

\textbf{Chollet \cite{Chollet2019}.}
Skill-acquisition efficiency, the core of the Abstraction and Reasoning
Corpus, is structurally close to $\eta_M$. Two differences: Chollet
treats prior knowledge $K$ as a confound to be controlled for in
measurement; here $K$ is constitutive, because $\eta_M$ operates on
$(A, K)$ and the efficiency of learning depends on what is already known.
Chollet's ``scope'' and ``generalization difficulty'' map to the growth
of the reachable territory $T_{\mathrm{reachable}}$ in the $(A, K)$
landscape.

\textbf{Schmidhuber \cite{Schmidhuber2010}.}
Curiosity as compression progress corresponds precisely to $\eta_M$:
the agent seeks states that maximize improvement in its world model. The
difference is foundational. Schmidhuber treats intrinsic drive as a
design objective, a reward signal to be engineered into artificial
agents. Theorem~\ref{thm:invariants}, Part~1, establishes $\gamma > 0$ as a
necessary condition of persistence: not a reward to be designed in, but a
property that selection enforces.

\textbf{Friston \cite{Friston2010, Friston2017}.}
The free energy principle derives minimization behavior from
self-organization: systems that persist must minimize surprise (free
energy). This is consistent with monotone descent on $H(W \mid K)$
\cite{McCarthy2026d}. The gap is the minimization \emph{imperative}:
FEP establishes that self-organizing systems \emph{do} minimize free
energy, but does not establish why $\gamma > 0$ is structurally necessary:
why entities that fail to minimize are eliminated. Theorem~\ref{thm:invariants},
Part~1, fills this gap via the lineage argument: $\gamma = 0$ lineages
contract under positive novelty rate $\nu(E) > 0$ regardless of their initial fitness.

\textbf{Cattell \cite{Cattell1971} and Spearman \cite{Spearman1904}: Psychometric intelligence.}
Cattell's distinction between fluid intelligence ($G_f$) and crystallized
intelligence ($G_c$) maps onto the framework. $G_f$, the capacity for
novel problem-solving independent of prior knowledge, is structurally
close to $\eta_M$: the rate at which the entity extracts grounding from
new learning signals. $G_c$, accumulated knowledge and skill, is
$|K|$, the size of the knowledge base. The framework predicts the
well-established correlation between $G_f$ and $G_c$: entities with
higher $\eta_M$ build larger $K$ faster, so the two measures co-vary
despite measuring different quantities.

Spearman's $g$-factor, the dominant factor across psychometric tests
, should, on this account, correlate with $I(E) = \gamma \cdot \eta_M$.
This is a testable prediction: $g$ should load on both engagement rate
(willingness to attempt novel problems) and learning efficiency
(improvement per attempt), and the product should predict $g$ better
than either factor alone.

% ------------------------------------------------------------------
\section{The Complete Chain}
\label{sec:chain}
% ------------------------------------------------------------------

The five papers derive the complete architecture of intelligence from a
single generative constraint. Each step is a consequence of bounded
active context interacting with properties of the physical world:

\begin{enumerate}
\item \textbf{Bounded $C_n$ + information preservation $\;\Rightarrow\;$
      graph structure} \cite{McCarthy2026a}. Any entity accumulating
      propositions under bounded context necessarily maintains a
      directed graph with typed edges.

\item \textbf{Bounded $C_n$ + growing $A$ $\;\Rightarrow\;$ convergence
      to escalation} \cite{McCarthy2026b}. Any viable control
      architecture converges toward escalation properties as concurrent
      demands grow relative to $C_n$.

\item \textbf{Survival pressure $\;\Rightarrow\;$ $K/A$ inseparability
      with asymmetric retention} \cite{McCarthy2026c}. Under survival
      pressure, knowledge and action become inseparable; the retention
      criterion is probability of need $\times$ conditional value.

\item \textbf{$\gamma > 0$ + feedback from $W$ $\;\Rightarrow\;$ gradient
      descent on $U$ with convergence; cost-of-failure $\;\Rightarrow\;$
      encoding hierarchy} \cite{McCarthy2026d}. Agency with feedback
      produces convergent uncertainty reduction. Asymmetric costs
      generate a permanent encoding hierarchy.

\item \textbf{Positive entropy + bounded $C_n$ $\;\Rightarrow\;$
      $\gamma > 0$ and $\eta_M > 0$ necessary; $I(E) = \gamma \cdot \eta_M$
      unique measure} [this paper]. The lineage argument and the
      context-filling argument establish the two invariants. Their
      product is the unique intelligence measure.
\end{enumerate}

Each step is a consequence of the same generative constraint. The
contribution is not any individual result but the connected chain: from
bounded context to graph structure to survival-driven retention to
escalation to evaluation-driven descent to the structural definition
of intelligence.

% ------------------------------------------------------------------
\section{Open Questions}
\label{sec:open}
% ------------------------------------------------------------------

\begin{enumerate}
\item \textbf{Origin of $\gamma$.} Theorem~\ref{thm:invariants} establishes
      that $\gamma > 0$ is necessary for persistence but does not explain how
      $\gamma > 0$ structures bootstrap from substrates where $\gamma = 0$. The
      thermodynamic conjecture: dissipative structures
      \cite{Prigogine1984, England2013} under sustained energy flow
      spontaneously produce $\gamma > 0$ configurations, which selection then
      retains. This is distinct from the necessity result.

\item \textbf{Direction of $\gamma$.} Is agency uniformly distributed across
      learning signals, or does it have intrinsic directedness? If
      directed, what determines the direction?

\item \textbf{The $M^{(n)}$ hierarchy.} How many levels of
      meta-learning does a persistent entity require? Is there a natural
      termination point, or is the hierarchy open upward without bound?

\item \textbf{Bound composition in $A$.} How do individual action
      bounds compose when multiple actions share knowledge
      substructure?
\end{enumerate}

% ------------------------------------------------------------------
\section{Conclusion}
\label{sec:conclusion}
% ------------------------------------------------------------------

Intelligence is the driven capacity to improve. This is not proposed but
identified by elimination. The lineage argument
(Theorem~\ref{thm:invariants}, Part~1) establishes that $\gamma = 0$ is
untenable: entities without agency contract under positive novelty rate
$\nu(E) > 0$ and are eliminated by selection. The
context-filling argument (Part~2) establishes that learning efficiency
$\eta_M > 0$ is necessary: without improvement, bounded context fills
with unresolved content and inference degrades to failure. Part~3
establishes minimality: every other property of a persisting entity can
equal zero for some persistent entity. The product $I(E) = \gamma \cdot
\eta_M$ is the unique measure consistent with these invariants,
derived from the law of total expectation applied to a probability gate
and a conditional rate.

The type irreducibility of $A$ and $M$
(Theorem~\ref{thm:type}) ensures that the distinction between executing
and learning is not an implementation choice but a categorical
necessity: no learning process can transform an action into a learning
process or vice versa.

The five-paper chain derives, from the single constraint of bounded
active context, the complete architecture: graph structure, knowledge--action
inseparability, escalation dynamics, evaluation-driven descent with
encoding permanence, and the structural invariants of intelligence. Each result
is a necessary consequence of the generative constraint interacting
with properties of the physical world: positive entropy, survival
pressure, and feedback. Nothing is stipulated; everything follows.

\section*{Acknowledgments}

The author used Claude (Anthropic) for drafting assistance, literature review, and \LaTeX{} preparation. All theoretical content, proofs, and arguments are the author's own.

\bibliographystyle{plain}
\begin{thebibliography}{99}

\bibitem{Ashby1956}
W.~R. Ashby.
\newblock \emph{An Introduction to Cybernetics}.
\newblock Chapman and Hall, 1956.

\bibitem{Barendregt1984}
H.~P. Barendregt.
\newblock \emph{The Lambda Calculus: Its Syntax and Semantics}.
\newblock North-Holland, revised edition, 1984.

\bibitem{Cattell1971}
R.~B. Cattell.
\newblock \emph{Abilities: Their Structure, Growth, and Action}.
\newblock Houghton Mifflin, 1971.

\bibitem{Chollet2019}
F.~Chollet.
\newblock On the measure of intelligence.
\newblock \emph{arXiv preprint arXiv:1911.01547}, 2019.

\bibitem{Cover2006}
T.~M. Cover and J.~A. Thomas.
\newblock \emph{Elements of Information Theory}.
\newblock Wiley-Interscience, 2nd edition, 2006.

\bibitem{DupouxLeCunMalik2026}
E.~Dupoux, Y.~LeCun, and J.~Malik.
\newblock Why {AI} systems don't learn and what to do about it: Lessons on
  autonomous learning from cognitive science.
\newblock \emph{arXiv preprint arXiv:2603.15381}, 2026.

\bibitem{England2013}
J.~L. England.
\newblock Statistical physics of self-replication.
\newblock \emph{The Journal of Chemical Physics}, 139(12):121923, 2013.

\bibitem{FinnEtAl2017}
C.~Finn, P.~Abbeel, and S.~Levine.
\newblock Model-agnostic meta-learning for fast adaptation of deep networks.
\newblock In \emph{Proceedings of the 34th International Conference on Machine
  Learning}, pages 1126--1135, 2017.

\bibitem{Friston2010}
K.~Friston.
\newblock The free-energy principle: A unified brain theory?
\newblock \emph{Nature Reviews Neuroscience}, 11(2):127--138, 2010.

\bibitem{Friston2017}
K.~Friston.
\newblock Active inference and learning.
\newblock \emph{Neuroscience and Biobehavioral Reviews}, 68:862--879, 2017.

\bibitem{Hutter2005}
M.~Hutter.
\newblock \emph{Universal Artificial Intelligence: Sequential Decisions Based
  on Algorithmic Probability}.
\newblock Springer, 2005.

\bibitem{LeggHutter2007}
S.~Legg and M.~Hutter.
\newblock Universal intelligence: A definition of machine intelligence.
\newblock \emph{Minds and Machines}, 17(4):391--444, 2007.

\bibitem{McCarthy2026a}
P.~D. McCarthy.
\newblock Graph structure is necessary for information preservation under
  bounded context.
\newblock Manuscript, 2026.

\bibitem{McCarthy2026b}
P.~D. McCarthy.
\newblock Convergence of control architectures to escalation under bounded
  context.
\newblock Manuscript, 2026.

\bibitem{McCarthy2026c}
P.~D. McCarthy.
\newblock Knowledge and action are inseparable under survival pressure.
\newblock Manuscript, 2026.

\bibitem{McCarthy2026d}
P.~D. McCarthy.
\newblock Evaluation-driven state revision and the encoding hierarchy
  under bounded context.
\newblock Manuscript, 2026.

\bibitem{Pearl1988}
J.~Pearl.
\newblock \emph{Probabilistic Reasoning in Intelligent Systems: Networks of
  Plausible Inference}.
\newblock Morgan Kaufmann, 1988.

\bibitem{Prigogine1984}
I.~Prigogine and I.~Stengers.
\newblock \emph{Order Out of Chaos: Man's New Dialogue with Nature}.
\newblock Bantam Books, 1984.

\bibitem{Schmidhuber2010}
J.~Schmidhuber.
\newblock Formal theory of creativity, fun, and intrinsic motivation
  (1990--2010).
\newblock \emph{IEEE Transactions on Autonomous Mental Development},
  2(3):230--247, 2010.

\bibitem{Schrodinger1944}
E.~Schr\"{o}dinger.
\newblock \emph{What Is Life? The Physical Aspect of the Living Cell}.
\newblock Cambridge University Press, 1944.

\bibitem{Shannon1948}
C.~E. Shannon.
\newblock A mathematical theory of communication.
\newblock \emph{Bell System Technical Journal}, 27(3):379--423, 1948.

\bibitem{Simon1955}
H.~A. Simon.
\newblock A behavioral model of rational choice.
\newblock \emph{The Quarterly Journal of Economics}, 69(1):99--118, 1955.

\bibitem{Spearman1904}
C.~Spearman.
\newblock ``General intelligence,'' objectively determined and measured.
\newblock \emph{The American Journal of Psychology}, 15(2):201--292, 1904.

\bibitem{Thrun1998}
S.~Thrun and T.~M. Mitchell.
\newblock Learning to learn.
\newblock In \emph{Learning to Learn}, pages 3--17. Springer, 1998.

\end{thebibliography}

\end{document}
